\chapter{المخرج الثاني}
% full version: chapters/16_chemoutput.tex

للآلة مخرجان. السريع هو العضلات. والبطيء هو كيمياء الجسم: الدماغ يتحكم في القلب والأوعية والغدد، ويرسل عبر الدم رسائل إلى الجسم كله دفعة واحدة. الدم أوسع نواقل البث بثًّا: تبلغ الرسالة الواحدة كل خلية في نحو دقيقة، ومعناها، كالعادة، يمنحه ”قفل“ المستقبِل.

\section*{دواسة الوقود والمكابح}

القلب ينبض وحده، وله ميقاتيته الخاصة. الدماغ لا يفعل إلا أن يعدّل الإيقاع بسلكين متضادين في الأثر. \nt{NORA} دواسة الوقود: تسرّع، لكن ببطء. و\nt{ACET} المكابح: تبطئ بسرعة، لأنها لا تحتاج إلى سلسلة طويلة من التفاعلات الكيميائية. وهنا أخيرًا نجد عملًا لنوع \nt{ADRE}: الأدرينالين المقذوف في الدم هو دواسة الوقود نفسها، لكنها مضغوطة للجسم كله دفعة واحدة.

\pencilfig{16}{\pencilart{16}}{يقود القلبَ سلكان متضادا الأثر: دواسة الوقود والمكابح.}

\section*{سلسلة بتغذية راجعة}

كثير من الهرمونات يُدار بسلسلة متعاقبة: يأمر الدماغ الغدة الأولى، وهي تأمر الثانية، والثانية تفرز الهرمون، والهرمون يخبر الدماغ أنه صار كافيًا. إنه ترموستات من ثلاث مراحل. ويعرف المهندس خطر مخطط كهذا: إذا جُعلت الحلقة ”حساسة“ أكثر مما ينبغي، بدأت تتأرجح، فلا يثبت المستوى بل يتذبذب. لذلك دقة هذا المنظم محدودة: لا يمكن جعله مستقرًا ودقيقًا جدًا في آن.

\pencilfig{127}{%
% three identical first-order stages with proportional feedback: secretion = max(0, K (target - level)).
% K = 3 settles at 3/4 of the target; K = 12 (above the stability limit K = 8) keeps oscillating. x = 0.42 t, y = 0.55 level
\foreach \yo/\t in {1.75/{حلقة ضعيفة: هادئة، لكن دون الهدف}, 0/{حلقة قوية: أقرب إلى الهدف، لكنها تتأرجح}}{
  \draw[pen] (0,\yo) -- (8.5,\yo);
  \draw[pcGray, densely dashed, line width=0.5pt] (0,\yo+0.55) -- (8.5,\yo+0.55);
  \node[lbl, anchor=south west] at (2.0,\yo+0.8) {\t};}
\node[lbl, anchor=west] at (8.55,2.3) {الهدف}; \node[lbl, anchor=west] at (8.55,0.55) {الهدف};
\begin{scope}[yshift=1.75cm]
\draw[penlite=pcBlue] plot[smooth] coordinates {(0.00,0.00) (0.17,0.01) (0.34,0.08) (0.50,0.19) (0.67,0.33) (0.84,0.46) (1.01,0.55) (1.18,0.59) (1.34,0.58) (1.51,0.55) (1.68,0.49) (1.85,0.43) (2.02,0.37) (2.18,0.34) (2.35,0.33) (2.52,0.34) (2.69,0.37) (2.86,0.40) (3.02,0.43) (3.19,0.44) (3.36,0.45) (3.53,0.45) (3.70,0.44) (3.86,0.42) (4.03,0.41) (4.20,0.40) (4.37,0.39) (4.54,0.39) (4.70,0.40) (4.87,0.40) (5.04,0.41) (5.38,0.42) (5.88,0.42) (6.38,0.41) (7.06,0.41) (8.40,0.41)};
\end{scope}
\draw[penlite=pcRed] plot[smooth] coordinates {(0.00,0.00) (0.17,0.05) (0.34,0.30) (0.50,0.68) (0.67,1.02) (0.84,1.23) (1.01,1.30) (1.18,1.26) (1.34,1.16) (1.51,1.02) (1.68,0.87) (1.85,0.72) (2.02,0.58) (2.18,0.47) (2.35,0.39) (2.52,0.38) (2.69,0.46) (2.86,0.57) (3.02,0.67) (3.19,0.71) (3.36,0.70) (3.53,0.65) (3.70,0.58) (3.86,0.50) (4.03,0.43) (4.20,0.41) (4.37,0.45) (4.54,0.53) (4.70,0.61) (4.87,0.65) (5.04,0.64) (5.21,0.60) (5.38,0.54) (5.54,0.47) (5.71,0.42) (5.88,0.42) (6.05,0.48) (6.22,0.56) (6.38,0.62) (6.55,0.64) (6.72,0.62) (6.89,0.56) (7.06,0.50) (7.22,0.44) (7.39,0.42) (7.56,0.45) (7.73,0.53) (7.90,0.60) (8.06,0.63) (8.23,0.62) (8.40,0.58)};
\node[lbl, anchor=north] at (4.2,-0.05) {الزمن};
}{نموذج لسلسلة هرمونية من ثلاث مراحل. التغذية الراجعة الضعيفة تهدّئ المستوى، لكنها تتركه دون الهدف بوضوح. والقوية تقربه أكثر، غير أن المستوى لم يعد يثبت، بل يتأرجح.}

\section*{نبضات لا مستوى}

بعض الهرمونات لا يعمل إلا بالنبضات. إذا أُعطيت باستمرار، ”تعبت“ الغدة المستقبِلة وتوقفت؛ وإذا أُعطيت بنبضات قصيرة مرة في الساعة، عملت. المستقبِل لا يقرأ المستوى، بل الإيقاع، كالمشبك في فصل شيفرة مورس الذي لا يسمع إلا التغيرات.

\section*{اليوم والليلة}

وأخيرًا، في الجسم مولّد بطيء دوره نحو يوم كامل. دوره الطبيعي يختلف قليلًا عن أربع وعشرين ساعة، وفي كل صباح يضبطه الضوء، كما يُضبط جهاز الراديو على المحطة. لا تخلط بينه وبين مولّد الساعة في الحاسوب: إنه لا يقيس خطوات الحساب. بل يبدّل أنماط عمل الآلة كلها: النهار والليل، اليقظة والنوم.

\fullversion{16}
